\documentclass[11pt]{article}
\usepackage[a4paper,top=2.0cm,bottom=2.0cm,left=2.2cm,right=2.2cm]{geometry}
\usepackage[T1]{fontenc}
\usepackage{lmodern}
\usepackage[english]{babel}
\usepackage{microtype}
\usepackage{amsmath,amssymb}
\usepackage{xcolor}
\usepackage[hidelinks]{hyperref}
\usepackage{enumitem}
\usepackage{parskip}
\setlength{\parskip}{0.4em}

% referee text vs our reply
\newenvironment{referee}
  {\begin{quote}\itshape\small}
  {\end{quote}}
\newcommand{\reply}{\par\noindent\textbf{Reply.}\ }
\newcommand{\changed}[1]{\textcolor{blue}{#1}}
\newcommand{\dd}{\mathrm{d}}

% Cross-reference numbers read out of paper2.aux; see
% provenance/make_crossrefs.py. Never type these by hand.
% ---------------------------------------------------------------
% GENERATED FILE -- do not edit by hand.
% Cross-reference numbers read out of paper2.aux by
% provenance/make_crossrefs.py, so that the response cannot cite a
% number the manuscript does not have.
% ---------------------------------------------------------------
\newcommand{\msSecDomain}{2.1}
\newcommand{\msSecBvp}{3.4}
\newcommand{\msSecClassification}{3.2}
\newcommand{\msSecInversion}{3.3}
\newcommand{\msSecThakurta}{2}
\newcommand{\msSecSubmersion}{2.2}
\newcommand{\msAppValidation}{Appendix A}
\newcommand{\msLemCompact}{2}
\newcommand{\msProtStatus}{1}
\newcommand{\msProtEndpoint}{2}
\newcommand{\msProtNames}{3}
\newcommand{\msProtEvidence}{4}
\newcommand{\msPropClassification}{4}
\newcommand{\msPropExteriorSep}{5}

\newcommand{\R}[1]{\par\medskip\noindent\textbf{#1}\par\nopagebreak}

\title{\vspace{-1.6cm}Response to the referees\\[2pt]
\large Manuscript CQG-116884\\[2pt]
\normalsize\emph{Brachistochrones in conformal Kerr spacetimes:
control domain, exterior separatrices, and adiabatic response}}
\author{Iman Rosignoli}
\date{\today}

\begin{document}
\maketitle
\thispagestyle{empty}

\section*{Summary of the revision}

I thank both referees for reports that were unusually specific and, in the first
case, unusually generous with technical detail. The first report identified a
defect that I had not seen and that turned out to be more serious than the
report could establish from outside; tracking it down produced a stronger result
than the one it replaced. I am grateful for it.

The revision follows the ``required reconstruction'' list of Referee~1 in the
order given there.

\begin{enumerate}[leftmargin=1.4em,itemsep=2pt]
\item \textbf{The exterior/contact/continuation distinction is now enforced
globally.} A new section~\msSecDomain, \emph{The control domain, and the status of
continued formulae}, proves as Lemma~\msLemCompact\ the compactness criterion the referee
sketched, derives from it that the compact-control problem \emph{ends} at
$r=2M$ for the selector $W=\partial_\eta$, and states as Protocol~\msProtStatus\ the
three-way labelling---(E) exterior extremal, (C) limiting contact, (A) analytic
continuation with no optimality claimed---which is then used without exception,
including in the figures.

\item \textbf{The extremal-versus-minimiser statements are corrected}, in both
papers, and the claim that numerical perturbation establishes minimality is
removed.

\item \textbf{The $J=-J_c$ cusp/corner contradiction is resolved by
rederivation}, not by editing a sentence. \emph{Both} printed statements were
wrong. The corrected result is a four-case classification
(Proposition~\msPropClassification), and all affected figures, captions, tables and conclusions
have been regenerated.

\item \textbf{The computational archive is rebuilt.} The missing module is
restored, the slopes and checksums are now \emph{generated} rather than
transcribed, and a manifest maps each artefact to one command and one checksum.
\end{enumerate}

The exterior retrograde separatrix has been promoted to the flagship exact
result, as recommended. The title no longer advertises a trichotomy, since there
is none.

Below, every comment is answered in order. Changes are marked in the highlighted
PDF.

\bigskip
\hrule
\section*{Referee 1}

\R{Major 1. The control domain ends at the conformal stationary limit.}

\begin{referee}
Timelikeness is therefore the reason the support maximum is finite and is
attained on a compact control set; it is not merely a convenient hypothesis.
[\dots] The current interior terminology must be changed globally, not only
qualified once in the abstract or conclusion.
\end{referee}

\reply Accepted in full; this is the structural change of the revision.

New section~\msSecDomain\ opens with the two rail constraints and proves
\textbf{Lemma~\msLemCompact}: the slice $\{g(u,u)=-1,\ -g(u,W)=\hat E,\ u$ future-directed$\}$
is nonempty and compact if and only if $W$ is timelike and $\hat E\ge|W|$,
degenerating to a single point exactly at $\hat E=|W|$, and is empty or
noncompact when $W$ is null or spacelike. The proof is the frame computation the
referee gave, which I have made explicit because the paper now leans on it.

Since $g(\partial_\eta,\partial_\eta)=-A^2(1-2M/r)$, the two scalar regularity
inequalities of Paper~I read here $r>2M$ and $\bar v^2=1-A^2f/\hat E^2>0$, so the
compact-control problem ends at $r=2M$. The text states explicitly that a change
of chart does not repair this---Doran coordinates regularise the formulae at
$\Delta=0$ but leave the scalar $g(W,W)$ unchanged---and that a Doran- or
ZAMO-anchored timelike selector would be a different problem, which is not
treated here.

Protocol~\msProtStatus\ fixes the labelling (E)/(C)/(A) and is applied throughout, including
in the figures: continued segments are drawn dashed and carry the legend entry
requested in M6.

\R{Major 2. Pontryagin extremals are not automatically local minima.}

\begin{referee}
It then incorrectly says that a Pontryagin extremal is ``guaranteed only locally
optimal.'' No such guarantee follows from the maximum principle alone.
\end{referee}

\reply The referee is right, and the error was worse than reported: the sentence
was present in \emph{both} versions of Paper~I, including one already under
consideration elsewhere. Both are corrected to state the hierarchy as the referee
writes it,
\[
  \text{minimiser}\Longrightarrow\text{PMP extremal},
  \qquad
  \text{PMP extremal}\not\Longrightarrow\text{local minimiser},
\]
with the explicit note that local minimality would require a second-variation or
conjugate-point argument, or another sufficient condition.

In Paper~II, ``controlled-rail extremal'' is now the default term; ``minimiser'',
``optimal'', ``fastest'' and the unqualified ``brachistochrone'' are reserved for
statements carrying a sufficient certificate over the same domain and endpoint
protocol. A remark in section~\msSecDomain\ states this rule.

The passage the referee will have had in mind opened Appendix~A:

\begin{referee}
The minimum principle is verified directly by perturbing the computed
brachistochrones and confirming that both $T_t$ and $T_\tau$ are minimized.
\end{referee}

\noindent It is gone. The appendix now says that such perturbations are a
consistency check on the numerics and cannot certify a minimum in an
infinite-dimensional admissible class.

\R{Major 3. Endpoint protocol and costate relabelling.}

\reply Accepted. Protocol~\msProtEndpoint\ fixes the default---fixed launch event, fixed
spatial target, free arrival clock, branch-dependent shooting---and states that
comparisons between the $\eta$- and $t$-branches are made at the relabelled
costate $J_t=A\,J_\eta$, so that ``same $J$'' and ``same endpoints'' are
different comparisons.

Section~\msSecBvp, the one place where the default does not apply, now names its two
protocols explicitly: (P1) same launch and (P2) two fixed spatial endpoints,
which give opposite depth orderings. The text records that in (P1) at a common
numerical $J$ the two branches compare \emph{different} orbits, which is exactly
why the protocols can disagree, and adds that an ordering between two selected
extremals does not show either is the global minimiser in its branch.

On the separate point: the statement that the printed Hamiltonians remain
algebraically regular at $f=0$ is now explicitly separated from the
control-domain question, at the end of section~2.

\R{Major 4. The retrograde marginal orbit: cusp or corner.}

\begin{referee}
This is a decisive internal contradiction. [\dots] I do not recommend resolving
the matter by choosing one sentence. The case $J=-J_c$ should be rederived from
the original Hamilton equations, including one-sided expansions, signs, branch
orientation, and numerical convergence tests.
\end{referee}

\reply Accepted, and this was the most valuable comment in the report. I
rederived the case as instructed. \textbf{Both} printed statements were wrong,
and the corrected result is sharper than either.

Two exact identities organise it. First,
$p_r^2=\mathcal S(r)/[D_E(r)^2\Delta(r)^2]$, the prefactor being regular and
nonzero at $r_e$, so the root structure of $p_r^2$ \emph{is} that of the shape
radical. Second, at the marginal value the turning bracket factorises exactly,
\[
  r\Delta-J_c^2D_E(r)=(r-2M)\big(r^2+J_c^2\big)
  \ \Longrightarrow\
  \mathcal S=r^4f^2w\,(r^2+J_c^2),
\]
a perfect square times a positive factor. Hence $r_e$ is a \emph{double} root for
\emph{both} signs of $J$, and the ``reality wall'' $\sqrt{wf}$ falls for both:
reality is not what selects the prograde value. The previous version's
counter-rotation argument is therefore withdrawn---and it was wrong in detail as
well as in role, since near $r_e$ the retrograde marginal has
$\dd\varphi/\dd r<0$ and $\dot r<0$, hence $\dot\varphi>0$, and is locally
\emph{co}-rotating.

What breaks the $\pm J$ symmetry is invisible to the shape radical, which
depends on $J$ only through $J^2$, and visible to the Hamilton equations:
$\tilde p_\varphi(r_e)=J-J_c$, which is linear in $J$. With
$\mathcal T=f/E-\tilde p_\varphi\varphi_0'$ the denominator of
$\dot r=R^2(\Delta/r^2)p_r/\mathcal T$, one has
$\mathcal T(r_e)=-a(J-J_c)/\bar P_e$, vanishing only at $J=+J_c$. The
classification is therefore fourfold, not threefold:

\begin{center}\small
\begin{tabular}{@{}llccl@{}}
\hline
case & $p_r^2$ at $r_e$ & $\mathcal T(r_e)$ & $r_e$ attained & shape\\
\hline
$|J|>J_c$ & $>0$ & $\neq0$ & no & smooth periapsis, exponent 2\\
$|J|<J_c$ & simple & $\neq0$ & finite parameter & semicubical cusp, exponent $3/2$\\
$J=-J_c$ & double & $\neq0$ & \textbf{asymptotic only} & finite tangent, exponent 1\\
$J=+J_c$ & double & simple & finite, and crosses & regular crossing\\
\hline
\end{tabular}
\end{center}

So the main text's ``cusp'' at $J=-J_c$ is wrong, and the appendix's ``corner''
is half right: the tangent is indeed finite, but there is \emph{no reflection}.
The parameter diverges logarithmically and $r_e$ is an asymptotic endpoint---the
behaviour a separatrix must have. The referee's caution against merging the
generic and marginal cases is now justified quantitatively: the cusp amplitude
\[
  K=\frac{EJ}{a^2}\sqrt{\frac{2M}{a^2-E^2J^2}}
\]
\emph{diverges} as $|J|\to J_c^-$, so the two families do not connect.

Closed forms are given for the retrograde decay rate,
$C=-\sqrt{4M^2+J_c^2}/(8M^2)$---which depends on $a$ and $E$ only through
$J_c$---for the prograde crossing rate, and for the one-sided tangent.

The one-sided expansions, signs, branch orientation and convergence tests the
referee asked for are in two scripts, \texttt{verify\_cusp\_corner.wls} (22
checks) and \texttt{verify\_marginal\_hamilton.wls} (23 checks), all passing.
Representative numbers at $M=1,a=0.9,E=1.2$: $C=-0.267000117$ against the closed
form to $10^{-6}$; the elapsed parameter gains $34.4956$ per four decades of
$r-r_e$ against $\ln(10^4)/|C|=34.4955$; the sub-marginal exponent measures
$0.4999999993$ against $1/2$.

All affected material has been regenerated, not recaptioned. The figures had
encoded the wrong law: the atlas drew a \emph{reflection} at $-J_c$, and the
classification bar painted the closed interval $[-J_c,+J_c]$ as cusp. The cusp
figure now carries the $-J_c$ curve and measures all three exponents, obtaining
$0.6668$, $1.0001$ and $1.9998$ against $2/3$, $1$ and $2$. Table~A1 is keyed on
root multiplicity and attainment rather than on a ``cusp bounce'' column that
could not hold $-J_c$. The abstract, introduction and conclusions are rewritten
accordingly, and the title no longer says ``trichotomy''.

\R{Major 5. The exterior separatrix should be the flagship exact result.}

\reply Accepted, and adopted as the organising result. Proposition~\msPropExteriorSep\ records
the exterior retrograde separatrix. Your values reproduce to sixteen digits from
the printed turning function:
\[
  r_d=3.513905124011657\,M,\qquad J_c^-=-8.053516003877019,
\]
a genuine double root, with the printed closed form of $J_c^-$ agreeing to
$10^{-12}$.

The proposition also certifies what makes it the right flagship: at $r_d$ both
control-domain inequalities hold strictly, $g(W,W)=-0.430833<0$ and
$\bar v^2=0.700811>0$, and they hold on the whole approach arc, so it is a
type-(E) controlled-rail extremal rather than a continuation.

One addition beyond the comment: a spin scan shows $r_d>2M$ persists, with
$r_d/M=2.857675,\,3.204756,\,3.513905,\,3.579494$ at $a=0.1,\,0.5,\,0.9,\,0.99$.
The exterior character is therefore structural rather than an accident of the
representative parameters, which strengthens the case for centring the paper on
it.

The prograde degeneration is treated exactly as you ask. A remark records
$r_d=1.512292\,M$, $J=2.936351$, notes that this lies outside the seed null
surface $r_+=1.435890\,M$ but inside $2M$, and states that it is an algebraic
degeneration of the analytically continued frozen equations---type~(A)---and not
a physical control separatrix for $W=\partial_\eta$. This also answers M4.

\R{Major 6. The fixed-endpoint result should retain its tiered proof status.}

\reply Accepted. The conclusions now keep the five levels apart explicitly
rather than merging them: (i) an analytic integrand inequality; (ii) a theorem
for $r_0\le R_*(E,a)$; (iii) an asymptotic result as $r_0\to\infty$ in the static
case; (iv) a computer-assisted proposition certified at representative $r_0$ over
a stated grid; (v) a conjecture in full generality. Protocol~\msProtEvidence\ (below) gives
these different evidential status by construction.

The endpoint protocol is fixed explicitly (major 3 above), and the text now says
that an ordering between two selected extremals does not by itself prove either
is the global minimiser in its branch.

On executability: you are right that the representative CAP runner imports
\texttt{cap\_full}, and right that the module was absent. See major 10.

\R{Major 7. Scope and numerical provenance of the first-order result.}

\begin{referee}
Figure 3 reports a common log--log slope $2.12\pm0.03$, whereas Table A3 reports
slopes from approximately 1.95 to 2.16. [\dots] the figure, table, raw data, fit,
uncertainty estimate, and file hash should be generated by the same archived
command.
\end{referee}

\reply Accepted, and the underlying problem was worse than the discrepancy you
saw: the numbers had been transcribed by hand from separate runs, and re-running
today gives $2.12\pm0.03$ and $2.15$, so neither printed value was current.
Re-copying them would only have reset the clock. They are now \emph{generated}.

A single archived command,
\texttt{python3 paper2/provenance/make\_provenance.py}, runs each validation
script as a subprocess, parses the slopes and their covariance $\sigma$, reads
the $\varepsilon$-window back out of the source rather than restating it, hashes
every script involved, and writes a machine-readable record, a \LaTeX{} fragment
of macros, and the manifest. The manuscript \texttt{\textbackslash input}s the
fragment, so the printed values cannot drift from the archive again. The raw
$(\varepsilon,\text{residual})$ pairs are dumped by the validation script, so
figure, table and any refit come from one dataset.

The theorem is now stated as local on a compact regular subarc, with turning
points, the freezing surface, the stationary limit, chart singularities and
moving-separatrix crossings excluded, and with passage through a slowly moving
separatrix identified as a singular inner problem rather than a uniform
continuation.

On the numbers themselves: you wrote that the values ``may all be compatible
with second-order residuals over a finite window''. That is now demonstrated
rather than asserted. Since the residual is $c_2\varepsilon^2(1+k\varepsilon+\cdots)$
the local exponent is $2+O(\varepsilon)$, so a finite-window fit is biased
upward, and the covariance $\sigma$ measures fit quality rather than that bias---%
which is why $2.12\pm0.03$ sits several $\sigma$ from $2$ without being in
tension with second order. Refitting on shrinking sub-windows gives
$2.118\to2.068\to2.043$ in every configuration, reaching $2.045$ at
$\varepsilon_{\max}=0.004$.

The $\varepsilon$-window is stated ($\{0.001,0.002,0.004,0.008,0.016\}$), as is
the fitted radial interval, $r/M\in[8,11]$, which remains $6M$ from the nearest
excluded locus.

\R{Major 8. Higher-genus terminology should remain neutral.}

\reply Accepted. The disclaimer was already present but was not binding on the
rest of the text. A convention now states that \emph{length-two iterated Abelian
integral on the genus-two hyperelliptic curve} is the only classification
claimed, and that ``genus-two dilogarithm'', wherever it appears, is an
explicitly descriptive shorthand for that object and not a statement that a
function class with specified kernels, relations, normalisation, monodromy or
single-valued completion has been established. The three theorem-level bare uses
have been converted.

``Elliptic dilogarithm'' is retained where it occurs, since at genus one that
object is established (Beilinson--Levin, Brown--Levin) and is not at issue.

\R{Major 9. Physical language should match the model and what is computed.}

\reply Accepted. Protocol~\msProtNames\ (Nomenclature) fixes the dictionary and is applied
throughout: $\hat E$ is the actively maintained rail charge and not the conserved
energy of a freely falling test particle; $J=p_\varphi$ is the axial control
costate and not mechanical angular momentum; the rail minimises a selected
elapsed clock and not thrust, fuel, impulse or integrated acceleration; $r_e=2M$
is the conformal stationary limit; $r_+$ is the seed Kerr null surface, or
conformal Killing-horizon candidate, with the explicit note that for
time-dependent $A(\eta)$ neither the apparent nor the event horizon has been
computed here; and inside $2M$ we write ``continued orbit'', ``root pattern'' or
``analytic branch'' rather than ``capture'', ``plunge'' or ``bounce''.

\R{Major 10. Reproducibility package.}

\begin{referee}
Some critical scripts are absent or are referenced in sibling directories outside
the archive; the CAP runners require the missing \texttt{cap\_full.py}; and
representative short file hashes in the cited release do not match the hashes
printed in the manuscripts.
\end{referee}

\reply All three findings are correct, and I am glad they were checked. Thank you
also for confirming that the exact reductions returned zero residuals.

\emph{The missing module.} \texttt{cap\_full.py} had been factored out of the CAP
prototype during development and never archived, so a fresh clone could not run
the representative computer-assisted proposition at all. It is restored,
extracted verbatim from the prototype so that the two cannot drift, with a
self-test. The runner now imports and certifies cells as expected.

\emph{Sibling directories.} The diagnosis turned out to be different from what it
looked like. All thirteen generators Paper~II needs were in fact tracked; but
every one of them inserted the repository root on \texttt{sys.path} while the
shared style module lives one directory down, so on a fresh clone every one
failed at import. Fixed in all thirteen. Fixing it also exposed a silent bug: a
turning-radius helper returned the grid point just inside the forbidden region,
so one curve was dropped from a log--log panel with no error raised.

\emph{Hashes.} The printed hashes were stale. In fact the same two files carried
three different values---one printed, one in the release you cloned, one in the
working tree---which is the signature of hand transcription. As described under
major 7, they are now generated by the same command that produces the slopes.

\emph{Manifest and environment.} \texttt{MANIFEST.tsv} maps sixteen artefacts,
each to one command and one checksum. The environment lockfile has been verified
against the live environment and extended with the computer-algebra system
version, which \texttt{pip} cannot record but the symbolic suite requires; the
lockfile is itself a manifest entry. A fresh clone no longer requires unarchived
parent or sibling folders.

\R{Major 11. Paper II depends on statements in Paper I that also require revision.}

\reply Accepted. The incorrect PMP sentence is corrected in Paper~I (major 2
above). Paper~I's existence and normality theorem has been expanded to define the
extended admissible class and clock state, state the regularity assumptions,
explain closure of the non-autonomous clock equation under the chosen
convergence, and identify the endpoint protocol and excluded degeneracies.

The Vaidya boundary qualification is carried into the abstract and conclusion:
the same compact-control domain ends at $r=2m(v)$, the calculation establishes an
exterior approach and contact threshold rather than an interior optimal
trajectory, and $m'(v)<0$ in the ingoing metric is a formal continuation, with
physical evaporation requiring the outgoing problem.

Paper~II no longer imports optimality claims from Paper~I; it cites the
verification criterion as a criterion.

\R{Major 12. Relation to established work.}

\reply Accepted, and treated as a positioning task as you suggest. A new
introduction paragraph separates what is established from what is new.
Established: the stationary problem and its reduction to a Finsler or Jacobi
length functional---Perlick's fixed-energy optical metric; the sub-Riemannian
formulation of Giannoni, Piccione and Verderesi; the arrival-time theory of
Giannoni and Piccione; and, for global minimality, the Morse theory of Giannoni,
Piccione and Tausk. I have added the Giannoni--Masiello--Piccione programme
(the timelike extension of Fermat's principle; the Morse theory covering massive
particles; and the finiteness results on \emph{conformally stationary}
spacetimes, which is the closest published setting to the frozen-$A$ regime used
here), and Caponio, Corona, Giamb\`o and Piccione on fixed-energy solutions of an
indefinite Lagrangian with an affine Noether charge, which is the closest
published variational structure to the controlled rail.

The distinction is stated explicitly at that point: in all of those the charge is
\emph{conserved}, because the generator is a symmetry; here it is \emph{held by
control} against a background of which $W$ is in general not a symmetry, at the
cost of proper acceleration, and $\hat E$ is a design parameter rather than an
integral of the motion. The three consequences with no stationary
counterpart---possible loss of compactness and hence a causal boundary, the
reduced Hamiltonian ceasing to be an interior first integral, and first-order
response as a genuine perturbation problem---are named there.

I have also taken up the framing of your closing paragraph: the emphasis is now
on the actively controlled invariant, the exact Hamiltonians, and the causal
boundary at which the compact control geometry changes category, rather than on
exhibiting another exact trajectory.

\R{Minor comments.}

\begin{description}[leftmargin=2.6em,style=nextline,itemsep=1pt]
\item[M1] Done: Protocol~\msProtEndpoint, stated once and used consistently, with the
distinct two-spatial-endpoint comparison flagged where it occurs.
\item[M2] Done: $J=p_\varphi$ is identified as a control costate at first use,
with the explicit statement that it is not mechanical angular momentum.
\item[M3] Done: the compact indicatrix boundary and the convex body whose support
function is used are now distinguished where the Hamiltonians are introduced,
with the reason no convexification is needed.
\item[M4] Done: see major 5. The interior double root is never called a physical
separatrix.
\item[M5] Done: units and the normalisation $M=1$ are declared once in
Protocol~\msProtNames and in every numerical caption and table; an audit of all captions
found one that listed parameters without $M$, now fixed.
\item[M6] Done: continued interior curves are drawn with a distinct line style
and carry the legend entry ``continuation, no optimality claimed''.
\item[M7] Done: Protocol~\msProtNames\ carries a table keeping $E$, $\hat E$, $E_{\rm eff}$,
$J$ and $J_{\rm eff}$ distinct, including the convention that $E$ abbreviates
$E_{\rm eff}$ in the frozen sections and the relabelling $J_t=A\,J_\eta$.
\item[M8] Done: plotted deviations are declared absolute unless labelled
relative, norms are named where the residual is a vector, and every fitted slope
is quoted with its interval and, for the adiabatic convergence, its
$\varepsilon$-window.
\item[M9] Done: Protocol~\msProtEvidence\ separates exact symbolic identities, independent
numerical agreement, recovery of established limits, and computer-assisted
propositions, and states their different evidential status. Claims are tagged.
\item[M10] Done: the conclusion is shortened and restructured, and lists only
results whose domain and proof status are established in the body.
\end{description}

\bigskip
\hrule
\section*{Referee 2}

\R{Bibliography.}

\reply Accepted with thanks; the list was useful and is now incorporated.
Perlick~(1991) was already central. I have added or made explicit: Giannoni,
Piccione and Verderesi (1991 sub-Riemannian formulation); Giannoni and Piccione
on arrival-time brachistochrones; Giannoni, Piccione and Tausk on Morse theory
for travel-time brachistochrones, with the \texttt{math-ph/9905007} preprint
identifier added so the 1999 preprint and the published version are visibly the
same work; Haws and Kiser (1995); and Ta\c{s} (\texttt{arXiv:2512.08776}).

Your suggestion that the global analysis and the Morse index be connected to that
literature is well taken and has shaped the revision: it is the framework in
which the referee-1 request to distinguish extremals from minimisers is properly
answered, and the introduction and Protocol~\msProtStatus\ now appeal to conjugate-point
and Morse-theoretic notions rather than asserting optimality.

One small bibliographic note, offered only so that the record is right: the
article at \emph{J. Math. Phys.} \textbf{38}(12), 6367--6381 (1997),
\texttt{doi:10.1063/1.532217}---the URL given as item~[6] in your list---is by
Giannoni, Piccione and Verderesi. Giannoni, Masiello and Piccione are a separate
and prolific collaboration, and I have now cited three of their papers in their
own right, including the finiteness results on conformally stationary spacetimes,
which are the closest published setting to the frozen-conformal-factor regime
used in section~\msSecClassification\ and which I had missed. Checking your list also caught an error
of mine: my own entry for that paper had the wrong page number, now corrected.

\R{Curvature-operator eigenvalues, weighted manifolds, and generalised
hyperplanes.}

\reply I have considered this carefully and, with respect, I do not think the
requested characterisation can be made, for two independent reasons.

First, the objects differ. A controlled-rail brachistochrone is not a geodesic of
a metric, weighted or otherwise: it is an extremal of a time-optimal control
problem with an actively enforced pointwise constraint $-g(u,W)=\hat E$, whose
admissible set is a compact indicatrix and whose Hamiltonian is a support
function. Its equations do not reduce to the geodesic equations of
$e^{-f^2(r)}g_{\mu\nu}$ for any $f$, because the constraint is not a conformal
rescaling of the metric; the revision makes this concrete by exhibiting a regime
where the control set loses compactness and the problem ceases to be posed at
all, which has no counterpart for geodesics of a weighted manifold.

Second, the requested statement---that the results here are ``a trivial
application'' of a theorem in a preprint---would be a claim about priority and
logical dependence that I am not in a position to make, since I do not find that
the cited theorem implies the controlled-rail construction, and since the works
in question are, so far as I can determine, unrefereed preprints. I would of
course reconsider on being shown the specific implication.

If the editor considers the connection worth recording, I would be glad to add a
neutral remark noting that weighted-manifold and curvature-operator methods
provide an alternative geometric setting for related variational problems,
without a dependence claim. I did not want to insert such a remark unilaterally.

\R{Experimental validation.}

\begin{referee}
The need is strict to present the techniques of experimental validation of the
findings.
\end{referee}

\reply I have addressed the substance of this, but must be candid about its
literal form: no experimental validation is available for this class of result,
and I do not claim any.

There are two independent reasons. Thakurta--Kerr is a conformal deformation of
Kerr whose central object is a compact object rather than a black hole, and for
time-dependent $A(\eta)$ the apparent and event horizons have not been computed;
and the controlled rail is an actively forced worldline requiring proper
acceleration, not a geodesic, so it is not the trajectory of any freely falling
body that is observed. Claiming observational contact would be precisely the
error that Referee~1 warns against in major 9.

What discharges the corresponding burden is now set out explicitly as
Protocol~\msProtEvidence, which distinguishes four levels of evidence: (I) exact symbolic
identities returning an exact zero in two independent computer-algebra systems;
(II) independent numerical agreement between genuinely different routes;
(III) recovery of established limits---$a\to0$ Schwarzschild, $A\to1$ the
stationary Kerr constants, $\hat E\to\infty$ the Fermat--Randers metric, and the
exact recovery of Perlick's fixed-energy optical metric; and (IV) computer-assisted
propositions with stated grids and windows. Level~(III) is the closest available
analogue to an independent experimental check, and it is falsifiable: a sign
error anywhere in the construction would break one of those limits. A remark
states plainly that no experimental level exists and why.

\bigskip
\hrule
\section*{A note of thanks}

Referee~1's report caused four errors of ours to be found: the cusp/corner
contradiction, the incorrect statement about Pontryagin extremals, the absent
\texttt{cap\_full.py}, and the drifting hashes. Pursuing them turned up three
more that the report could not have seen from outside---a broken import path that
prevented every figure script from running on a fresh clone, a silent
turning-radius bug that dropped a curve from a published panel, and a wrong page
number in a reference. The manuscript is materially better for it.

\end{document}
